\setcounter{section}{0}

\chapter*{Approfondir son sujet}
\subsection*{\centering Objectif : Passer de la recherche à l'argumentation solide.}

\vspace{0.5cm}

\hrule
\vspace{0.2cm}

\textbf{Question Définitive Retenue :} \dotfill

\vspace{0.2cm}
\hrule
\vspace{0.5cm}

\section{La Chasse aux Exemples et Arguments (Le Convaincant)}
\begin{enumerate}[label=\textbf{A.\arabic*}]
    \item \textbf{L'Exemple Phare (L'Accroche) :}
    \small Décrivez l'anecdote factuelle ou l'analogie la plus parlante que vous utiliserez pour introduire votre sujet et capter l'attention du jury.\\
    \normalsize
    \framebox[\linewidth]{\parbox{\linewidth}{\vspace{0.3cm}\dotfill \vspace{0.3cm} \dotfill \vspace{0.3cm}}}

    \item \textbf{Arguments Forts (Axe 1 : Thèse, Causes, Avantages) :}
    \small Identifiez trois arguments principaux pour votre développement. Chacun doit être étayé par un fait.
    \normalsize
    \begin{itemize}[label=--]
        \item Argument 1 (Idée principale) : \dotfill \quad \small (Exemple ou fait associé : \dotfill)
        \item Argument 2 (Idée principale) : \dotfill \quad \small (Exemple ou fait associé : \dotfill)
        \item Argument 3 (Idée principale) : \dotfill \quad \small (Exemple ou fait associé : \dotfill)
    \end{itemize}

    \item \textbf{Arguments Limites (Axe 2 : Antithèse, Conséquences, Solutions) :}
    \small Identifiez trois arguments qui nuancent ou proposent une solution à la problématique.
    \normalsize
    \begin{itemize}[label=--]
        \item Argument 1 (Idée principale) : \dotfill \quad \small (Solution ou limite associée : \dotfill)
        \item Argument 2 (Idée principale) : \dotfill \quad \small (Solution ou limite associée : \dotfill)
        \item Argument 3 (Idée principale) : \dotfill \quad \small (Solution ou limite associée : \dotfill)
    \end{itemize}
\end{enumerate}

\vspace{0.5cm}
\hrule
\vspace{0.5cm}

\section{Analyse des Sources et Maîtrise du Vocabulaire}

\begin{enumerate}[label=\textbf{B.\arabic*}]
    \item \textbf{Maîtrise du Vocabulaire Technique (Clarté pour le Jury) :}
    \small Remplissez le tableau avec les termes clés de votre sujet, puis proposez une analogie ou une explication simple pour un non-spécialiste (ex: professeur de philosophie).
    \normalsize
    \begin{center}
    \begin{tabular}{|>{\centering\arraybackslash}m{3cm}|m{4cm}|m{7.5cm}|}
    \hline
    \textbf{Mot-clé NSI} & \textbf{Définition Brève} & \textbf{Analogie / Explication Simple} \\
    \hline
    \small Exemple : \emph{Oversampling} & \small Technique de réplication de données pour équilibrer la base. & \small Quand il manque des données d'un ingrédient, on le réplique pour que la recette ne soit pas biaisée. \\
    \hline
    \dotfill & \dotfill & \dotfill \\
    \hline
    \dotfill & \dotfill & \dotfill \\
    \hline
    \dotfill & \dotfill & \dotfill \\
    \hline
    \end{tabular}
    \end{center}

    \item \textbf{L'Analyse Critique des Sources (Rigueur) :}
    \small Identifiez et analysez vos deux meilleures sources.
    \normalsize
    \begin{itemize}[label=--]
        \item \textbf{Source Experte (Rapport, Livre, Publication) :} Titre et Auteur : \dotfill
        \item Apport à mon sujet (Notions NSI validées) : \dotfill
        \item \textbf{Source Vulgarisée (Vidéo YT réputée, Article de Presse de Qualité) :} Titre et Auteur : \dotfill
        \item Apport à mon sujet (Exemples, chiffres percutants) : \dotfill
    \end{itemize}

\end{enumerate}